% ============================================================
% 分片：s06b-models  对应原文 PDF 第 82--86 页及第 87 页顶部（印刷页码同）
% 负责章节：6.2 Some Models
% 译者：t05-s32   审校状态：初稿
% 审校：t06-s33 已完成逐式核对（修正 0 处）
% ============================================================
% 边界说明：本分片自 PDF 第 82 页 §6.2 标题起，至第 87 页顶部 §6.2 收尾（跨页句、
% "We have, for illustrative purposes…" 整段、致谢段）结束；第 87 页 "REFERENCES"
% 标题及其后文献 1--53 归 s07-references，本分片不取。
% [边界: 第87页 REFERENCES 之前的 §6.2 收尾与致谢属本分片；参考文献归 s07]
% 原文脚注（第 82 页 ⋆）以 \footnote{} 保留，编号由 LaTeX 重排。
% 原文用直立体 d 表示微分，故微分记号统一写作 \mathrm{d}。

% ==== 原文 PDF 第 82 页（印刷页 82）====
\label{op:82}
\subsection{一些模型}

现在将这套机制应用于某些具有 $(0,2)$ 超对称的 $\sigma$ 模型。为简单起见，我们只考虑靶空间为 $\mathbf{CP}^{n-1}$ 中的超曲面 $X$ 的情形。

我们从 §3 已经知道，构造一个在低能下流向 $X$ 的 $\sigma$ 模型的超可重整 $(2,2)$ 模型需要哪些超场。为本文目的，必须把这些超场分解为 $(0,2)$ 多重态，以便最终构造 §3 模型的 $(0,2)$ 形变——以及它们的 Landau--Ginzburg 描述。

于是考虑 $(0,2)$ 超空间中的 $U(1)$ 规范理论，其中含有如下 $(0,2)$ 手征多重态：$n$ 个电荷为 1 的场 $S_i$，\footnote{此处所有超多重态均为 $(0,2)$ 超多重态，记号中省略撇号。}一个电荷为 $-n$ 的场 $P$，以及一个中性的场 $\Sigma$（来自 $(2,2)$ 规范多重态的约化）。此外，还需要 $n$ 个电荷为 1 的 Fermi 多重态 $\Lambda_{-,i}$，其
\begin{equation}\tag{6.35}
E_i=i\sqrt2\Sigma S_i,
\end{equation}
以及另一个电荷为 $-n$ 的 Fermi 多重态 $\Lambda_{-,0}$，其
\begin{equation}\tag{6.36}
E_0=-in\sqrt2\Sigma P.
\end{equation}
为完全确定该模型，还需选取额外的函数 $J^i$（$i=1\ldots n$）与 $J^0$，满足
\begin{equation}\tag{6.37}
E_iJ^i+E_0J^0=0.
\end{equation}
在 $(2,2)$ 情形下，我们曾取超势 $W=PG(S_i)$，其中 $G$ 是齐次的 $n$ 次多项式，并满足通常的横截性条件。用 $(0,2)$ 语言，

% [Tutor Note: (6.37) 不是任意的技术条件：§6.1 中 Fermi 多重态由 $D_+\Lambda_{-,a}=\sqrt2E_a$ 定义，而超对称要求为其配备满足 $E_aJ^a=0$ 的 $J^a$（见 (6.33)--(6.34) 前后的讨论）。可用本模型的具体表达式验证：$E_iJ^i+E_0J^0=i\sqrt2\Sigma S_i\cdot P\frac{\partial G}{\partial S_i}-in\sqrt2\Sigma P\cdot G=i\sqrt2\Sigma P\left(S_i\frac{\partial G}{\partial S_i}-nG\right)=0$，其中用了欧拉定理 $S_i\partial G/\partial S_i=nG$。(6.36) 的负号正是使两项抵消之所需；若把 $E_0$ 的符号取反，(6.37) 便不成立。]

% ==== 原文 PDF 第 83 页（印刷页 83）====
\label{op:83}
% [跨页接续: 原文第82页/第83页]
按照 (6.34)，这一选取给出
\begin{equation}\tag{6.38}
\begin{aligned}
J^i&=P\frac{\partial G}{\partial S_i}\\
J^0&=G.
\end{aligned}
\end{equation}
为得到 $(0,2)$ 模型，我们将在保持 (6.37) 及规范不变性的前提下改变这些选取。我们还将假定 $J^0$ 仍与 $P$ 无关，而 $J^i$ 仍线性于 $P$。（规范不变性所允许的其他项在手征超场中将是更高阶的，因而在重整化群意义下看来是“无关”的。对 $E$ 的可能修改也是如此。）因此不妨保留 $J^0=G$ 作为 $G$ 的定义。但 $J^i$ 是可以改变的。一般可能性是
\begin{equation}\tag{6.39}
\begin{aligned}
J^i&=P\frac{\partial G}{\partial S_i}+PG^i\\
J^0&=G,
\end{aligned}
\end{equation}
其中 $G^i$ 是 $S_k$ 的 $n-1$ 次多项式，选取得使
\begin{equation}\tag{6.40}
S_iG^i=0.
\end{equation}
于是得到一族 $(0,2)$ 模型，除通常的数据外还由 $G^i$ 参数化。把相关公式汇总起来，理论的势能为
\begin{equation}\tag{6.41}
\begin{aligned}
U(s_i,p)={}&\frac{e^2}{2}\left(\sum_i|s_i|^2-n|p|^2-r\right)^2+|G|^2\\
&+|p|^2\sum_i\left|\frac{\partial G}{\partial s_i}+G^i\right|^2+2|\sigma|^2\left(\sum_i|s_i|^2+n^2|p|^2\right).
\end{aligned}
\end{equation}
这里 $S_i$、$P$ 与 $\Sigma$ 的玻色分量依次记作 $s_i$、$p$ 与 $\sigma$。因此可以像 $(2,2)$ 情形那样分析对 $r$ 的依赖。对 $r\gg0$ 得到低能理论为靶空间 $X$ 的 $(0,2)$ $\sigma$ 模型的相。对 $r\ll0$ 得到 $(0,2)$ Landau--Ginzburg 相。我们已把 C-Y/L-G 对应推广到 $(0,2)$ 模型。下面更仔细地考察这两个相。

% [Tutor Note: (6.41) 可逐项对上 §6.1 的 (6.27)：$U=\frac{e^2}{2}(\sum_iQ_i|\phi_i|^2-r)^2+\sum_a|E_a|^2+\sum_a|J^a|^2$。第一项是消去 $D$ 后的 $D$ 项势，$D=e^2(\sum_i|s_i|^2-n|p|^2-r)$，其中 $P$ 的荷 $-n$ 给出 $-n|p|^2$；$|E_i|^2+|E_0|^2=2|\sigma|^2(\sum_i|s_i|^2+n^2|p|^2)$ 正是末项（$S_i$ 荷为 1 故系数 2，$P$ 荷为 $-n$ 故系数 $2n^2$）；$|J^0|^2=|G|^2$ 与 $\sum_i|J^i|^2=|p|^2\sum_i|\partial G/\partial s_i+G^i|^2$ 给出中间两项。核对 $\sigma$ 质量项时应先确认这一荷的平方关系。]

% ==== 原文 PDF 第 84 页（印刷页 84）====
\label{op:84}
\subsubsection*{Landau--Ginzburg 相}

对 $r\ll0$，$p$ 获得期望值，与 $(2,2)$ 情形一样，规范群破缺为 $\mathbb{Z}_n$。因此该模型是一个 $\mathbb{Z}_n$ 轨形。

无质量多重态是手征多重态 $S_i$ 与 $\Lambda_{-,i}$；其他多重态都有质量。低能理论中 $E_i$ 与 $J^i$ 的有效值通过把 $\sigma$ 与 $p$ 取为它们的真空期望值——即 0 与 $\sqrt{-r/n}$——得到。于是在有效的 $(0,2)$ Landau--Ginzburg 模型中
\begin{equation}\tag{6.42}
\begin{aligned}
E_i&=0\\
J^i&=\langle p\rangle\left(\frac{\partial G}{\partial s_i}+G^i\right).
\end{aligned}
\end{equation}
当然，这些满足 $E_iJ^i=0$。

% [Tutor Note: 在 $r\ll0$ 的相中，$\langle p\rangle=\sqrt{-r/n}\neq0$ 而 $\langle\sigma\rangle=0$：前者由 $D=0$ 定出，后者因为 $\sigma$ 从末项质量项 $2|\sigma|^2(\sum_i|s_i|^2+n^2|p|^2)$ 获得质量（在 $s_i=0$、$|p|^2=-r/n$ 处质量平方为 $-2nr>0$）。于是 $E_i=i\sqrt2\Sigma S_i$ 的有效值因 $\langle\sigma\rangle=0$ 而为零，$J^i$ 只剩下 $\langle p\rangle(\partial G/\partial s_i+G^i)$ 的部分，$E_iJ^i=0$ 自动成立。]

\subsubsection*{Calabi--Yau 相}

对 $r\gg0$，无质量的玻色模式是与 $X$ 相切的模式。右行无质量费米子正是无质量玻色子的 $(0,2)$ 伙伴。

但左行无质量费米子如何？这里的情形更有趣。$\Lambda_{-,0}$ 的费米分量有质量，但 $\lambda_{-,i}$（超场 $\Lambda_{-,i}$ 的费米分量）的某些模式是无质量的。通过检查上述拉格朗日量公式可以看出，无质量模式满足
\begin{equation}\tag{6.43}
\sum_i\bar s_i\lambda_{-,i}=0
\end{equation}
以及
\begin{equation}\tag{6.44}
\sum_i\left(\frac{\partial G}{\partial s_i}+G^i\right)\lambda_{-,i}=0.
\end{equation}
下面把这些公式纳入 $(0,2)$ $\sigma$ 模型的一般框架。一般而言，无质量左行费米子是靶空间 $X$ 上某个全纯矢量丛 $F$ 的截面。这意味着必须能够用

% [Tutor Note: (6.43) 与 (6.44) 分别给出 $n$ 个 $\lambda_{-,i}$ 的两个线性约束，其中 (6.43) 对 $s_i$ 不是全纯的（含 $\bar s_i$），(6.44) 则是全纯的。两个独立约束把 $\lambda_{-,i}$ 的 $n$ 维空间截到 $n-2$ 维，这正是靶空间维数 $\dim_{\mathbb{C}}X=n-2$，也提示 $F$ 应是一个秩 $n-2$ 的丛；当 $G^i=0$ 时它应还原为 $TX$（见下页）。$P$ 的费米伙伴 $\lambda_{-,0}$ 有质量，是因为 $\langle p\rangle\neq0$ 使 $\Lambda_{-,0}$ 与 $\Lambda_{-,i}$ 通过 $E$、$J$ 数据耦合。]

% ==== 原文 PDF 第 85 页（印刷页 85）====
\label{op:85}
% [跨页接续: 原文第84页/第85页]
在 $s_i$ 中全纯变化的方程来描述无质量模式。现在，(6.44) 对 $s_i$ 是全纯的，而 (6.43) 不是。为获得更多理解，应把 (6.43) 替换为一个有相同后果、但对 $s$ 全纯变化的方程。这不难做到。代替约束 (6.43)，引入一个规范不变性，规定 $\lambda_i$ 的两个模式若由变换
\begin{equation}\tag{6.45}
\lambda_i\to\lambda_i+s_i\lambda
\end{equation}
相联系则视为等价，其中 $\lambda$ 为某个量。于是 (6.43) 可看作一个规范固定条件；不变性 (6.45) 的每条轨道都有唯一满足 (6.43) 的代表。

为把 (6.45) 与 (6.44) 放入它们的理论背景，引入 $\mathbf{CP}^{n-1}$ 上通常的线丛 $\mathcal{O}(k)$（限制到 $X$ 上），并考虑 $X$ 上矢量丛的一个映射序列
\begin{equation}\tag{6.46}
0\to\mathcal{O}\xrightarrow{\ \alpha\ }\bigoplus_{i=1}^{n}\mathcal{O}(1)\xrightarrow{\ \beta\ }\mathcal{O}(n)\to0,
\end{equation}
其中 $\alpha$ 是映射 $\lambda\to s_i\lambda$，而 $\beta$ 是映射
\begin{equation}\tag{6.47}
\lambda_i\to\lambda_i\left(\frac{\partial G}{\partial s_i}+G^i\right).
\end{equation}
(6.46) 是矢量丛的一个复形（即 $\beta\alpha=0$），但它不是正合序列。相反，$\ker\beta/\operatorname{im}\alpha$ 是 $X$ 上秩为 $n-2$ 的全纯矢量丛 $F$。

% [Tutor Note: $\beta\alpha=0$ 直接可验：$\beta(\alpha(\lambda))=\beta(s_i\lambda)=s_i\lambda(\partial G/\partial s_i+G^i)=\lambda\left(s_i\partial G/\partial s_i+s_iG^i\right)=n\lambda G$（用了欧拉定理与 (6.40)），而 $X$ 恰由 $G=0$ 定义，故在 $X$ 上此式为零。这也说明为何要先把线丛限制到 $X$ 上：复形性质只在 $X$ 上成立。]

当 $G^i=0$ 时，$F$ 约化为 $X$ 的切丛 $TX$。一般地，$F$ 是 $TX$ 的一个形变；反过来，$TX$ 作为 $X$ 上全纯矢量丛的所有形变都可以用 (6.46) 描述。这一 $TX$ 的形变族此前已被讨论过 [52, 第 II 卷, pp. 461-3]。$F$ 的选取（连同其余数据）确定一个以 $X$ 为靶空间的 $(0,2)$ 模型，因此 $G^i$ 参数化这样一族模型；我们已找到了该族的 Landau--Ginzburg 描述。

% [Tutor Note: (6.46) 给出的是一种“扩张型”形变：$F=\ker\beta/\operatorname{im}\alpha$ 的秩为 $n-2$，与 $\dim X$ 相符；$G^i=0$ 时 $\beta$ 与 $\alpha$ 的核/像给出通常的欧拉序列 $0\to\mathcal{O}\to\oplus_i\mathcal{O}(1)\to TX\to0$（对 $\mathbf{CP}^{n-1}$ 中的超曲面），故 $F=TX$。非零 $G^i$ 相当于在 $\beta$ 中加了一个保持 (6.40) 与 $\beta\alpha=0$ 的扰动，从而把 $TX$ 形变为另一个全纯矢量丛。]

% ==== 原文 PDF 第 86 页（印刷页 86）====
\label{op:86}
人们可能会问：这一通常 C-Y/L-G 对应的 $(0,2)$ 类比能得出什么具体结果。一个回答（沿 §3.3 的思路）是：$\operatorname{Tr}(-1)^F$ 与椭圆亏格在从 Calabi--Yau 到 Landau--Ginzburg 模型的跃迁下不变；而对半扭 $(0,2)$ 模型（它决定低能 Yukawa 耦合）而言，Landau--Ginzburg 是 Calabi--Yau 的解析延拓。（A 型与 B 型拓扑场论没有 $(0,2)$ 类比。）另一个回答（沿 §5.5 的思路）是：涉及 $(0,2)$ 模型的拓扑改变过程可以像 $(2,2)$ 模型一样分析。

为把故事讲完，或许还应指出：在左行无质量费米子取值于靶空间 $X$ 上某个矢量丛 $F$ 的 $(0,2)$ $\sigma$ 模型中，对象 $E_a$ 与 $J^a$ 在几何上应分别理解为 $F$ 与 $F^*$（$F$ 的对偶）的全纯截面。若它们非零，这些截面会给某些模式以质量。在我们一直研究的这个具体例子中（$X$ 为 $\mathbf{CP}^{n-1}$ 中的超曲面、$F$ 如上），有效的 $E_a$ 与 $J^a$ 为 0。它们几乎不可能是别的，因为在这些例子中 $F$ 与 $F^*$ 没有非零的整体全纯截面。但在 $(0,2)$ $\sigma$ 模型的一般研究中，当 $F$ 使得 $F$ 或 $F^*$ 具有整体截面时，$E_a$ 与 $J^a$ 的考虑很可能极为重要。

% [Tutor Note: “$E_a$ 是 $F$ 的截面、$J^a$ 是 $F^*$ 的截面”意味着 (6.37) $E_aJ^a=0$ 是 $F\otimes F^*\to\mathcal{O}$ 的配对给出的一个全纯函数（在 $X$ 上应为常数、在这里为零），这正是 $E$、$J$ 数据的几何含义；在本节的超曲面例子里 $F=\ker\beta/\operatorname{im}\alpha$ 与 $F^*$ 都没有非零整体全纯截面，所以有效 $E_a$、$J^a$ 只能为零，Landau--Ginzburg 相中 $\langle\sigma\rangle=0$ 给出的 $E_i=0$ 与此一致。]

\subsubsection*{共形不变性？}

在 §3 中，我们构造 $(2,2)$ 模型的 C-Y/L-G 对应时并不需要知道共形不变性在两侧是否成立。在本节中，我们对 $(0,2)$ 模型做了同样的事。

对 $(0,2)$ 模型，情况被认为与 $(2,2)$ 模型相当不同。§3 中所研究的 $(2,2)$ 模型都被认为具有共形不变的不动点，并在红外流向它们。在微扰论逐阶意义上，上面分析的 $(0,2)$ 模型也是如此。但在许多情形下，这些近似的 $(0,2)$ 共形不动点会被非微扰修正破坏 [46,53]。例如，对于我们一直仔细考察的这个具体情形，据信当
% [跨页接续: 原文第86页/第87页]
% ==== 原文 PDF 第 87 页（印刷页 87）====
\label{op:87}
多项式 $G$、$G^i$ 取一般值时就会发生这种情况。从世界面的观点看，人们对共形不变性破坏的本质知之甚少。
% [原文疑误: 原文 "Little is understand, from the world-sheet point of view, about …" 中 "understand" 应为 "understood"（形容词性过去分词），系笔误；译文按 "understood" 理解。]

为作说明，我们只详细研究了可以作为 $(2,2)$ 模型扰动来构造的 $(0,2)$ 模型。用本文的技巧研究其他 $(0,2)$ 模型将会非常有趣。

致谢：作者感谢 P. Aspinwall、B. Greene 与 D. Morrison 的有益讨论，以及他们在早期阶段介绍自己的工作。

% [Tutor Note: 这是 Witten 对 $(0,2)$ 情形下 C-Y/L-G 对应是否成立的判断：在微扰论层面，（线性 $\sigma$ 模型两侧的）对应与 $(2,2)$ 一样成立；但 $(0,2)$ 模型未必有真正的红外共形不动点——非微扰效应可以破坏近似的固定点，故不能像 $(2,2)$ 那样断言两侧各自存在一个共形场论再谈它们等价。这与本节开头“不预先假定共形不变性”的构造策略相呼应：对应关系是在线性 $\sigma$ 模型的层面上建立的，而非在 CFT 层面上。]
